% Owner: classification_e. Root and boundary_d full-source referee PASS.
\begin{proposition}[E.2.18 (Proved): the normalized cusp and its marked hexagon]
Let $\nu:\mathcal S\to\mathcal W$ be the
simultaneous normalization of the relative cusp
in E.2.17, and let $\mathcal D\subset\mathcal S$
be its reduced conductor.  After shrinking the
parameter disc, each fibre pair is the standard
toric pair
\[
 (\mathcal S_s,\mathcal D_s)
 \simeq(\mathrm{dP}_6,\hbox{its toric hexagon}).
\]
The six conductor components may be labelled over
the disc so that their incidence, self-intersections
$-1$, opposite-pair involution, and endpoint
identifications are those of the punctured family.
The relative logarithmic tangent bundle has a
global holomorphic frame:
\[
 T_{\mathcal S/\Delta}(-\log\mathcal D)
 \simeq\mathcal O_{\mathcal S}^{\oplus2}.
\]
This assertion identifies the normalized pair
and supplies a relative frame; it does not by
itself identify the analytic smoothing neighbourhood.
\end{proposition}
\begin{proof}
In the relative charts of E.2.17, the normalized
branches are smooth over the original parameter.
Their conductor branches and their intersections
are also smooth over that parameter.  The
normalization of the conductor is consequently
a proper smooth family of curves, and the
normalized intersection strata are finite
etale over the disc.  Their connected components
and incidence are constant.  Since the disc
is simply connected, label the six curve
families $\mathcal D_1,\ldots,\mathcal D_6$
and the six vertex sections by continuation
from a punctured fibre.  Each curve fibre
is a smooth curve of genus zero.  Its two
distinct vertex sections remain distinct;
the same relative local charts exclude a
new self-intersection or an additional
incidence at the centre.  Thus the central
conductor is a hexagon of six embedded
rational curves.  The degree of each
relative normal line is constant in its
proper curve family, so all six central
self-intersections equal $-1$.

Write $S=\mathcal S_0$ and $D_i=\mathcal D_{i,0}$.
Ehresmann transport gives $b_1(S)=0$ and
$b_2(S)=4$.  Contract the three disjoint
curves $D_1,D_3,D_5$.  The analytic
Castelnuovo contraction gives a smooth
compact complex surface $Y$; the needed
analytic statement, including smoothness
of the resulting point, is
\citet[\S4, no.~8(a), pp.~364--365]{Grauert1962E}.
Only these three contractions on the
central surface are used; no relative
contraction theorem is assumed.

The images $C_2,C_4,C_6$ of the remaining
curves are smooth rational curves.  Each
meets two contracted curves transversely
once, so has self-intersection
$-1+1+1=1$.  Each pair meets once, at
the image of their intervening contracted
curve; these three intersection points
are distinct and the intersections are
transverse.  Thus they form a triangle.
Also $b_1(Y)=0$ and $b_2(Y)=1$.

Any one of these curves, say $C=C_2$,
has nonnegative intersection with every
effective divisor: distinct irreducible
curves have nonnegative local intersection
multiplicities and $C^2=1$.  Adjunction
gives $K_Y\cdot C=-3$.  A nonzero section
of $K_Y^{\otimes r}$ for $r>0$ would
give an effective divisor of intersection
$-3r$ with $C$, impossible.  Hence
$\kappa(Y)=-\infty$.  The compact-surface
classification
\cite[Chapter~VI, Theorem~1.1 and Table~10]{BHPV2004E}
together with $b_1(Y)=0$ makes $Y$ rational:
class~VII has first Betti number one,
and a ruled surface has the first Betti
number of its base curve.  A rational
surface with $b_2=1$ is $\mathbb P^2$.
The curves $C_2,C_4,C_6$ are lines,
since their self-intersections are one.
Their three pairwise intersection points
are distinct, so they form a projective
triangle.  Blowing up its three vertices
recovers exactly the toric degree-six
del Pezzo surface and its hexagonal
boundary.  This proves the assertion
about the central pair; the punctured
pairs were already the prescribed ones.

The conductor involution is holomorphic
over the disc.  Its permutation of the
six labelled curve families, and its
permutation of the vertex sections,
are locally constant.  They are therefore
the same opposite-pair and endpoint
identifications as on a punctured fibre.
The remaining scalar gluing parameters
are not claimed to be determined by
this discrete observation.

Finally put
$\mathcal E=T_{\mathcal S/\Delta}(-\log\mathcal D)$.
For every fibre pair just identified,
the two toric infinitesimal actions
trivialize its logarithmic tangent
bundle.  Thus
\[
 \mathcal E|_{\mathcal S_s}\simeq
 \mathcal O_{\mathcal S_s}^{\oplus2},
 \qquad h^0(\mathcal S_s,\mathcal E_s)=2.
\]
Constant-cohomology base change, in the
form used in E.3.1, makes the direct
image a rank-two holomorphic vector
bundle on the disc and identifies its
fibres with these section spaces.
Choose a frame after shrinking the
disc.  The evaluation map from its
pullback to $\mathcal E$ is an
isomorphism on every fibre, hence
is an isomorphism everywhere by
Nakayama's lemma.  This gives the
claimed relative frame without
choosing a potentially singular
family of projective coordinate
changes on the punctured disc.
\end{proof}
